\cleardoublepage

% additional.tex
\section{Additional Concepts / \kr{추가 개념}}

This section covers essential vocabulary and techniques that will help you master the fundamental concepts of numbers and their relationships. These additional tools complement the main concepts and provide practical methods for solving problems.

\subsection{Place Value Vocabulary / \kr{자리값 용어}}

Understanding place value requires familiarity with specific terminology. The essential place values you need to learn include billions, hundred millions, ten millions, millions, hundred thousands, ten thousands, thousands, hundreds, tens, and ones.

\blockquote{\kr{한국어의 자리값 체계와 비슷하게, 왼쪽으로 갈수록 10배씩 커진다. 예: 일, 십, 백, 천, 만, 십만, 백만, 천만, 억, 십억, 백억}}

\blockquote{\kr{ones(일), tens(십), hundreds(백), thousands(천), ten thousands(만) 영어는 천부터 x10}}

For decimal place values, the key terms are tenths, hundredths, thousandths, ten-thousandths, and millionths.

\blockquote{\kr{소수의 자리값은 뒤에 ‘th’를 붙여 나타낸다. 예: tenths(소수 첫째 자리), hundredths(소수 둘째 자리), thousandths(소수 셋째 자리)}}

\blockquote{\kr{반드시 알아야하는건 tenths(할), hundredths(푼), thousandths(리)}}

\subsection{Fraction Fundamentals}

Fractions represent parts of a whole, and understanding their components is crucial. Every fraction has two main parts: the numerator and the denominator.

\blockquote{\kr{분수(fraction)는 분자(numerator)와 분모(denominator)로 이루어져 있다.}}

\blockquote{\kr{외우기 힘들다면 AND라는 줄임말 기억하면 된다. AND는 분수 = 분자/분모(numerator/denominator)의 순서를 의미한다.}}

There are different types of fractions that serve different purposes. An improper fraction has a numerator larger than its denominator, while a mixed number combines a whole number with a proper fraction.

\blockquote{\kr{가분수(improper fraction)는 분자가 분모보다 큰 분수이다. 대분수(mixed number)는 자연수와 분수가 결합된 형태이다.}}

\blockquote{\kr{이는 성질만 이해하면 쉽게 외울 수 있다. 가분수는 분자가 분모보다 크므로 ‘부적절한(improper)’ 형태이고, 대분수는 자연수와 분수가 혼합된(mixed) 형태이다.}}

\subsection{Decimal Types / \kr{소수의 종류}}

Decimals come in different forms, each with unique characteristics. A repeating or recurring decimal has a pattern that continues infinitely, while a finite decimal has a limited number of digits after the decimal point. An infinite decimal continues without end.

\blockquote{\kr{순환소수(repeating/recurring)는 일정한 패턴이 끝없이 반복되는 소수이다. 유한소수(finite)는 소수 자리가 일정 부분에서 끝나는 소수이다. 무한소수(infinite)는 끝없이 이어지는 소수이다.}}

\blockquote{\kr{repeating/recurring(반복), in-(없다) + finite(제한이 있는) = infinite(무제한)}}

\blockquote{\kr{순환소수의 구별방법은 소수점위에 점의 유무 예: \(1.\overline{3}\) = 1.333… 와 같이 위에 점으로 표시한다.}}
\blockquote{\kr{예: \(1.\overline347\overline5\)에서 3과 5 위에 점이 찍혀 있다면, 그 구간이 순환되어 \(1.347534753475\ldots\) 이 된다.}}
 

\blockquote{\kr{무한소수는 분수로의 변환이 불가하다}}

\blockquote{\kr{순환소수는 분수로 변환할 수 있으며, 순환되는 자리수를 9로 나누어 표현한다. 예: \(1.\overline{3984}\)의 경우 \(3984/9999 = 0.398439843984\ldots\) 로 나타낼 수 있다.}}


\newpage
\subsection{Comparison Vocabulary / \kr{부등식 용어}}

Mathematics uses specific language to describe relationships between numbers. Understanding these terms helps you interpret problems correctly and communicate mathematical ideas clearly.

For greater than relationships, we use terms like greater than, more than, larger than, longer than, higher than, above, exceeds, and over. For greater than or equal to, we say greater than or equal to, at least, more than or equal to, minimum, no less than, not under/below, or not smaller than.

For less than relationships, we use less than, fewer than, smaller than, shorter than, lower than, under, below, or beneath. For less than or equal to, we say at most, does not exceed, maximum, no more than, up to, or not above/greater than.

\blockquote{\kr{특히 at least(최소)와 at most(최대)는 반드시 알아두어야 한다. sat 수학에서 자주 사용된다.}}

\blockquote{\kr{무작정 외우기보다는 문제를 읽을 때 의미를 함께 생각해 보면 이해하기 쉽다.}}

\newpage
\subsection{Prime Factorization / \kr{소인수분해}}

Prime factorization is the process of breaking down a number into its smallest building blocks--prime factors. This technique is fundamental for understanding number structure and solving problems involving factors and multiples.

\blockquote{prime factorization(\kr{소인수분해}) \kr{는 말 그대로 수를 소수(prime number)로 완전히 분해하는 것을 의미한다.}}

\blockquote{\kr{소인수분해의 기본 순서는 먼저 어떤 수를 분해할지 확인하는 것이다.}}

\blockquote{\kr{그 수가 어떤 배수인지 확인한 후 소수로 나누어 간다.}}

To determine if a number is divisible by certain factors, we can use these rules:
\begin{itemize}
  \item \kr{2의 배수: 짝수/일의 자리 숫자가 2, 4, 6, 8, 0인 수}
  \item \kr{3의 배수: 각 자리의 수의 합이 3의 배수}
  \item \kr{끝의 두 자리 수가 00이거나 4의 배수}
  \item \kr{5의 배수: 일의 자리 숫자가 0또는 5}
  \item \kr{7의 배수: 세자리수 abc일 때 ab-c*2가 7의 배수인수 (자주 사용되지는 않는다)}
\end{itemize}

\blockquote{\kr{위의 규칙으로 배수를 판단한 후, 작은 소수부터 차례대로 나누어 소인수분해를 진행한다.}}

\subsection{Finding GCF and LCM}

The greatest common factor (GCF) and least common multiple (LCM) are essential concepts for working with multiple numbers. While these ideas may seem complex, there's a simple method using the Korean "\kr{ㄴ}" technique that makes finding them straightforward.

\blockquote{\kr{한국 학교에서는 ‘ㄴ’자 방법을 사용한다.}}

\blockquote{\kr{한국 학교에서 이미 배웠겠지만, ‘ㄴ’자 방법을 사용하면 쉽게 구할 수 있다.}}

Let's find the GCF of 12 and 18 using this method:

\[
\begin{array}{r|cc}
6 & 12 & 18 \\
\hline
  & 2 & 3
\end{array}
\]

\blockquote{\kr{마지막에 남은 숫자가 모두 소수가 되면 소인수분해가 완료된다.}}

\blockquote{\kr{최대공약수(GCF)는 ‘ㄴ’자 옆에 있는 숫자들의 곱으로 구한다. 예: 6.}}

Now let's find the LCM of 12 and 18:

\[
\begin{array}{r|cc}
6 & 12 & 18 \\
\hline
  & 2 & 3
\end{array}
\]

\blockquote{\kr{최소공배수(LCM)는 ‘ㄴ’자 옆과 밑에 있는 모든 숫자들의 곱으로 구한다. 예: $6 \times 2 \times 3 = 36$.}}


These methods provide efficient ways to work with factors and multiples, making complex problems more manageable and helping you develop a deeper understanding of number relationships.
